\subsection{Implemented Components}
\label{sec:mechanisms}
We implemented an experiment controller that starts isolated and concurrent
jobs, schedules request arrivals independently and enforces the CPU allocation
for each process. It checks thread affinity during execution, saves timing
and resource samples, and collects outputs before removing its experiment pod.

The controller connects inference to a document worker that reuses unchanged
embeddings, publishes index versions atomically and retrieves evidence for the
new request. Separate checks record answer agreement, citation validity and
the retrieved version. The framework also saves predictions and their hashes
before the final timing test, along with the inputs, configuration and source
hashes needed to trace each result. These are the components we built around
llama.cpp, MiniLM and existing data and scoring tools. We do not claim that
these general capabilities were absent from all previous systems.

\paragraph{Earlier Distributed Precision Control}
A separate four-node pilot changes only stage 0 of Qwen2.5-14B between existing
Q4\_K\_M/Q3\_K\_M artifacts. At a token boundary it exports KV, destroys and
unmaps the old executor, loads the target, imports state and checks re-export
integrity. Replaying the full prefix through all stages checks target-model accuracy. This validates
a transition path rather than a new quantizer or selective repair algorithm.
The detailed lifecycle is retained in the separate artifact appendix.

\paragraph{CPU Execution and Fixed Attention Omission}
The later monolithic worker keeps the full Q4\_K\_M model loaded, distinguishes
readiness from request timing and clears KV before each request. Four threads
serve prefill and two serve decode, with an explicit prompt boundary so a
singleton prefill tail remains in the prefill phase. Each logical frame contains 64 tokens. The physical microbatch size controls
how the graph is split. Runtime repacking
changes the CPU weight layout for existing quantized kernels \citep{llamacpp}.
It neither changes the GGUF file nor introduces a new quantization method.

We test raw microbatches of 2, 8, 32 and 64, and repacked microbatches of 2 and 64. Larger batches and
layout changes are opt-in approximate execution configurations: the raw-mb2
worker reproduces reference logits at 64/565 tokens, but that equality is not
assumed for other configurations. Repacking can increase startup and anonymous
memory, so those costs are measured separately from warm requests.

The fixed mask omits attention sublayers 39--42, including their Q/K/V/O
operations and KV writes, while retaining every FFN. The mask is constant
for a session. Online reactivation is not implemented: missing KV and changed
downstream states prevent treating it as a free reversible action. Attention
dropping is established prior work \citep{attentiondrop}. Residual similarity suggests candidate layers to omit. Answer checks determine
whether a configuration passes.

\paragraph{What Each Change Tests}
Precision changes target memory use. Batching and repacking target execution time.
Attention omission removes work at potential quality cost. Update deferral
protects the current request by delaying freshness. Each is evaluated at its
own endpoint before being connected to utility. A successful implementation,
a smaller process counter and an acceptable application outcome are different
claims. We implement no controller that jointly selects these actions online.
